% Created 2026-05-11 Mon 12:30
% Intended LaTeX compiler: pdflatex
\documentclass[11pt]{article}

\usepackage[utf8]{inputenc}
\usepackage[T1]{fontenc}
\usepackage{graphicx}
\usepackage{longtable}
\usepackage{wrapfig}
\usepackage{rotating}
\usepackage[normalem]{ulem}
\usepackage{amsmath}
\usepackage{amssymb}
\usepackage{capt-of}
\usepackage{hyperref}
\usepackage{minted}
\usepackage{tabularx}
\author{flavio}
\date{\today}
\title{introduzione a java}
\hypersetup{
 pdfauthor={flavio},
 pdftitle={introduzione a java},
 pdfkeywords={},
 pdfsubject={},
 pdfcreator={},
 pdflang={English}}
\begin{document}

\maketitle
\tableofcontents

\section{Tipi di dato base}
\label{sec:org1ef7355}
I tipi di base sono built-in, non sono creati dall'utente. Serve esserne a conoscenza per ragion di efficenza.
\subsection{Un tipo di dato descrive un insieme di valori e di operazioni definite su quei valori:}
\label{sec:org1e9179f}
\begin{table}[htbp]
\label{Tipi di dato}
\centering
\begin{tabular}{lll}
Tipo & Dominio & Operatori\\
\hline
int & Interi & +-*/\%\\
double & Numeri in virgola mobile & \\
boolean & valori booleani & \&\&\\
char & caratteri & + -\\
\hline
 &  & \\
\end{tabular}
\end{table}
\subsection{Tutti i tipi built-in}
\label{sec:org55a36cb}
\begin{center}
\begin{tabular}{lll}
Tipo & Intervallo & Dimensione\\
\hline
float & parte intera: ±1038, parte frazionaria: circa 7 cifre  decimali significative & 4 byte\\
double & parte intera ±10308, parte frazionario: circa 15 cifre decimali significative & 8 byte\\
int & -2147483648…2147483647 & 4 byte\\
long & -2147483648…2147483647 & 4 byte\\
boolean & true,false & 1 byte\\
char & rappresenta tutti i caratteri codificati con unicode & 2 byte\\
short & -32768…32767 & 2 byte\\
byte & -128…127 & 1 byte\\
\hline
 &  & \\
\end{tabular}
\end{center}
\subsection{Precedenza operatori aritmetici}
\label{sec:orgdbf817d}
\begin{center}
\begin{tabular}{lll}
Operatori & Operazioni & Precedenza\\
\hline
*,/,\% & Moltiplicazione,divisione,resto & Valutati per primi, da sinistra verso destra\\
+,- & Addizione,Sottrazione & Valutati per secondi, da sinistra verso destra\\
\hline
 &  & \\
\end{tabular}
\end{center}
\section{Variabili, dichiarazioni e assegnazioni}
\label{sec:org3686f22}
Una variabile è un nome usato per riferirsi a un valore di un tipo di dato
\begin{itemize}
\item Può essere dichiarata (\texttt{int a;}) e quindi non è detto che abbia un valore, se non in certi contesti specifici.
\begin{itemize}
\item Il valore viene assegnato a una variabile mediante assegnazione \texttt{a = 0;}
\begin{itemize}
\item Si può anche dichiarare e assegnare allo stesso tempo: \texttt{int a = 0;}

\begin{minted}[]{java}
    int a,b; // dichiarazione
    a = 5; // assegnazione
    b = a+10; // altra assegnazione
    int c = a+b; // dichiarazione e assegnazione
    a = c-3; // assegnazione
\end{minted}
\end{itemize}
\end{itemize}
\end{itemize}
\section{Letterali}
\label{sec:org57f578f}

Un letterale è una rappresentazione nel codice sorgente del valore di un tipo di dato.

\begin{itemize}
\item 27 e 32 sono letterali per gli \texttt{interi}
\begin{itemize}
\item 3.14 è un letterale per i \texttt{double}
\begin{itemize}
\item \texttt{true} o \texttt{false} sono gli unici due letterali per il tipo booleano
\begin{itemize}
\item \texttt{"ciao mondo"} è un letterale per il tipo String
\end{itemize}
\end{itemize}
\end{itemize}
\end{itemize}
\subsection{per distinguere gli interi da i numeri in virgola mobile}
\label{sec:org02a2479}
\begin{itemize}
\item Le costanti di tipo \texttt{int} sono semplicemente numeri nell'intervallo [-2000000000,+2000000000]
\begin{itemize}
\item Le costanti di tipo \texttt{long} sono specificate con il suffisso \texttt{l} o \texttt{L}
\end{itemize}
\end{itemize}
\subsection{Virgola mobile}
\label{sec:orge819e1b}
\begin{itemize}
\item \texttt{double}: numeri con la virgola
\begin{itemize}
\item \texttt{float}: hanno il suffisso \texttt{f} o \texttt{F}: \texttt{10.5f}
\end{itemize}
\end{itemize}
\section{Espressioni}
\label{sec:orgb103d47}
Un' Espressione è un letterale, una variabile o una sequenza di operazioni su letterali o varibili che producono un valore.

\begin{minted}[]{java}
int a = 42*(x-2)+y;

\end{minted}
\section{Caratteri e Stringhe}
\label{sec:orgb3852ea}
Un \texttt{char} è un carattere alfanumerico o un simbolo. Ci sono \(2^{16}\) possibili valori di caratteri.
\subsection{Caratteri di escape}
\label{sec:org35e8fde}
In alcuni casi dobbiamo rappresentare caratteri speciali o non stampabili.
\subsubsection{breve lista}
\label{sec:orgb8f9b73}
\begin{itemize}
\item Tab: \texttt{'\textbackslash{}t'}
\item A capo: \texttt{'\textbackslash{}n'}
\item Backlash: \texttt{'\textbackslash{}\textbackslash{}'}
\item Apice: \texttt{‘\textbackslash{}’’}
\item Virgolette: \texttt{‘\textbackslash{}”’}
\end{itemize}
\section{Priorità operatori}
\label{sec:org14cc343}
\begin{center}
\begin{tabular}{ll}
Operatori & Operazioni\\
\hline
Post-incremento e post-decremento & ++var --var +espr -espr\\
Operatori moltiplicativi & * / \%\\
Operatori additivi & + -\\
shift & << >> >>>\\
relazionali & < > <= >= instanceof\\
confronto & \textasciitilde{} \texttt{= !} \textasciitilde{}\\
And bit a bit & \&\\
XOR & \^{}\\
OR & \\
OR logico & \\
operatore ternario & ?:\\
lambda & ->\\
assegnazione & = += -= *= /= \%= \&= \^{}= <<= >>= >>>=\\
\end{tabular}
\end{center}
\section{Conversioni di tipo}
\label{sec:orgb04ae99}

Supponiamo di voler utilizzare il programma per effettuare dei calcoli. Come convertiamo da \texttt{stringa} a \texttt{intero}?
\subsection{conversione automatica a stringa}
\label{sec:org5eb016c}
Quando cerchiamo di concatenare a una stringa un operando di tipo diverso viene convertito automaticamente a stringa.

\begin{minted}[]{java}
public class PensieroProfondo
    {
        public static void main(String[] args)
            {
                String s = "La risposta alla fondamentale sulla vita, l'universo e tutto quanto e' ";
                int v = 42;
                String risposta = s+v; // equivalente a String risposta = s+Integer.parseInt(v)
                System.out.println(risposta);
            }
}
\end{minted}
\subsection{Conversione esplicita}
\label{sec:orgf23f2e1}
Utilizzando un metodo che prende in ingresso un argomento di un tipo e restituisce un valore di un altro tipo.
\texttt{Integer.parseInt(int n), Double.parseDouble(double n), Math.round(double n),  Math.floor(double n), Math.ceil(double n)}
\subsection{Cast esplicito}
\label{sec:org4dc7f91}
Anteponendo il tipo desiderato tra parentesi all'espressione essa verrà convertita.
\begin{minted}[]{java}
System.out.println((int)2.781828);
\end{minted}
\subsection{Cast implicito}
\label{sec:orgce3815f}
\subsubsection{Avviene o in fase di assegnazione: byte, short e char possono essere prommossi a int.}
\label{sec:orgc006094}
-int può essere promosso a long
-float a double
\subsubsection{calcolo espression}
\label{sec:orgab2f96b}
se uno dei due operandi è double, l'intera espressione è promossa a double. Altrimenti se uno dei due operandi è float, l'intera espressione è promossa a float
\section{Istruzioni di controllo}
\label{sec:org3ad6d93}
Finora abbiamno controllato il flussso di esecuzione attraverso La sequenza (l'ordine di esecuzione delle istruzioni). Mediante le istruzioni di controllo condizionali possiamo controllare il flusso del programma.
\subsection{if ed else}
\label{sec:org084de92}
Per realizzare una decisione si usa l'istruzion \texttt{if}.
\begin{itemize}
\item La sintassi è \texttt{if(<espressione booleana>) solo un istruzione}
oppure:
\texttt{if (<espressione boolean>)} \{
\}

\begin{minted}[]{java}
if (x < 0) x = -x; // calcola valore assoluto
if(x < y){ //credit default swap
    int t = x;
    x = y;
    y = t;

}
\end{minted}
\item l'\texttt{else} \textbf{che è facoltativo} viene eseguito se la condizione è falsa. Bisogna fare attenzione quando si annidano i diversi if ed else: l'else si riferisce sempre all'ultimo .
\begin{minted}[]{java}
if(x > 0)
    if(y > 0)
        System.out.println("x e y sono  > 0");
else
    System.out.println("x <= 0")
if(x > 0)
    if(y > 0)
        if(z < 5)
            if(w >= 3)
                System.out.println("x e y >= 0 e z <= 5 e w >= 3")
\end{minted}
\end{itemize}
\subsection{Switch}
\label{sec:org5425616}
Per confrontare una espressione intera si può usare l'espressione \texttt{switch}
\begin{minted}[]{java}
String hello = null;
switch(new Random.nextInt(6))
    {
        case 0: hello = "Ciao" break;
        case 1: hello = "hello" break;
            // etc ...

    }
//versione più compatta
String s = switch(c) // converte un carettere in rappresentazione
{
    switch 'k' -> "kappa";
    switch 'h' -> "acca";
    switch 'l' -> "elle";
    default -> "non so leggerlo";
}
\end{minted}
\section{Istruzioni iterative}
\label{sec:org38f18ec}
Ci sono 3 tipi di istruzioni iterative: \texttt{while, for, do while}
\subsection{while}
\label{sec:org48b87ea}
sintassi:
while(<espressione>)
\{
fino a quando espressione è vera
\}
semantica:
eseguito fino
\subsection{break e i cicli annidati}
\label{sec:orgd89fb56}
Il break ci permette di uscire solamente dal ciclo corrente se non si specifica una etichetta.
\begin{minted}[]{java}
outer:
for (int i = 0; i < h; i++)
    {
        for (int j = 0; j < w; j++)
            {
                if (i == j) break outer;
            }
    }


\end{minted}
\end{document}
